\section{Empirical Strategy}
\label{sec:empirical_strategy}

\subsection{Institutional Setting}

Before Morris, final land-use review in New York City ran through the Board of Estimate.
The Board combined three citywide officials, who together held six votes, with the five borough presidents, each holding one vote.
As the Supreme Court's opinion describes, the Board exercised broad fiscal and legislative powers, including authority over land use, franchises, contracts, and parts of the budget process \citep{board_of_estimate_of_city_of_new_york_v_morris_notitle_1989}.
The key implication is that land-use decisions were made by a body whose smallest geographic representative was the borough, a large swath of the city, rather than a district-based councilmember accountable to a much narrower constituency.

The 1989 Charter revision induced by the Morris decision substantially reorganized that authority. After the Board was abolished, final review of many land-use actions shifted to the City Council \citep{lane_policy_1998}.
The Council is a 51-member districted legislature.
Applications still pass through the community board, borough president, and City Planning Commission, but the Council now holds the final legislative stage for many discretionary approvals.
While the mayor can veto some Council actions, the practical politics of district-specific land use are dominated by the Council.

The key informal institution in this setting is member deference: the practice by which the Council generally follows the local councilmember on land-use applications within that member's district.
The 2025 Charter Revision Commission describes this practice as operating outside the formal Charter but central to how land use actually works \citep{new_york_city_charter_revision_commission_charter_2025}.

The post-Morris shift did not necessarily imply an immediate local-member veto. 
The 1989 Charter revision created the institutional conditions for district-level land-use control by moving many final land-use decisions into a 51-member Council. But the informal norm of member deference appears to have hardened over time. 
During the 1990s, Speaker control could still override local members; by the 2000s and 2010s, the local member’s position had become increasingly decisive. This delayed institutional consolidation is important for the empirical design: the expected effect is not an immediate 1990 break, but a growing homeowner gradient as member deference became binding.

\subsection{Unit of Analysis}

In this first Council-district version, I use the 51 City Council districts from the 2010 redistricting cycle as the unit of analysis.
This is the political geography through which member deference operates, and it is the relevant geography for a post-2010 descriptive test.
The cost is that the geography is not available consistently for the full pre-period, so I treat the 2010 districts as fixed and project historical exposure into those boundaries.

This is deliberately a first pass.
I do not claim that the 2010 districts were the operative political districts in the 1980s or 1990s.
Instead, the exercise asks whether the places that became homeowner-heavy 2010 Council districts, measured using 1990 tenure, are the places that lost large-building production once member deference became more important.

Housing production is measured from a single construction source: 25v4 MapPLUTO.
I assign each MapPLUTO lot to the archived 2010 Council district containing the lot's representative point, then aggregate \texttt{yearbuilt} and residential units to Council-district-by-year cells.
This keeps the construction outcome source fixed over 1980--2025.
The tradeoff is that MapPLUTO provides a surviving-stock year-built proxy, not a clean annual completion-flow series.

\subsection{Treatment Definition}

The treatment captures baseline homeowner concentration relative to the borough.
I start with 1990 NHGIS tract owner-occupied and occupied housing counts, intersect the 1990 tract polygons with the 2010 Council-district boundaries, and allocate tract counts to Council districts by area share.
Let $H_{db,1990}$ be the resulting 1990 owner-occupied share of occupied housing units in Council district $d$, assigned to majority borough $b$, and let $H_{b,1990}$ be the corresponding borough-wide owner-occupied share.
I define the treatment as a within-borough standardized homeownership score:

\begin{equation}
    T_{db}
    =
    \frac{
        H_{db,1990} - \overline{H}_{b,1990}^{\mathrm{Council}}
    }{
        s_b(H_{db,1990})
    },
\end{equation}

where $\overline{H}_{b,1990}^{\mathrm{Council}}$ and $s_b(H_{db,1990})$ are the mean and
standard deviation of 1990 homeownership across 2010 Council districts in
borough $b$.

This score compares each Council district to others assigned to the same borough. A value of zero indicates average homeownership for the borough; positive values indicate a more homeowner-heavy district; negative values indicate a more renter-heavy one.

This within-borough standardization matters because boroughs differ substantially in housing stock, density, transit access, baseline ownership rates, and long-run development pressure.
A citywide measure would partly conflate Staten Island with Manhattan or outer Queens with the Bronx.
The within-borough score instead asks whether homeowner-heavy Council districts within the same borough experienced different post-reform construction trajectories than their lower-homeownership neighbors.

The identifying logic is straightforward. If member deference gives local councilmembers effective control over discretionary land use, then places with stronger homeowner constituencies should have greater political capacity to resist large buildings.
That prediction should be sharpest for multi-unit and large-building construction, where local opposition and discretionary review are most consequential, and weaker for small buildings, which are less likely to require the same political process.
Figure~\ref{fig:treatment_map} illustrates the variation in the treatment measure across the city, with substantial within-borough dispersion that permits comparisons free of borough-level confounds.

\begin{figure}[htbp]
    \centering
    \includegraphics[page=1,width=0.82\textwidth]{../tasks/build_ccd2010_homeownership_1990_measure/output/ccdist2010_homeownership_1990_map.pdf}
    \caption{2010 Council-District Homeownership Exposure}
    \label{fig:treatment_map}
\figurenotes{Figure maps the first-pass treatment measure by 2010 City Council district. The measure is the 1990 homeownership share, built by area-weighting 1990 NHGIS tract tenure counts into the 2010 Council boundaries and standardizing within assigned majority borough. Redder districts had higher homeownership than other Council districts in the same borough; bluer districts had lower homeownership.}
\end{figure}
